\chapter{Quantum Algorithms and Cryptanalysis}
\label{ch:quantum-algorithms}

\abstract*{This chapter provides a precise treatment of the two quantum algorithms that reshape the cryptographic landscape: Shor's algorithm, which renders integer factorization and discrete logarithms tractable, and Grover's algorithm, which provides a quadratic speedup for unstructured search. We analyze the quantum resources required to break deployed cryptographic parameters and establish the boundary between quantum-vulnerable and quantum-resistant problems.}

% =============================================================
% SOURCE: notes.tex §3.1 (lines 1870–2501)
% REUSE: ~70%. Add formal circuit descriptions, resource
%        estimate tables, QRAM discussion.
% =============================================================

\section{The Quantum Computing Model}
\label{sec:qc-model}

% TODO: Brief recap (readers are assumed to know QM)
% - Qubits, gates, circuits, measurement
% - Quantum Fourier transform
% - The circuit model vs adiabatic/measurement-based
% - Focus on what's needed for Shor and Grover

\section{Shor's Algorithm}
\label{sec:shor}

% TODO: Adapt from notes.tex §3.1.1
% - Period finding as the core subroutine
% - Reduction of factoring to period finding
% - Quantum circuit: Hadamard, modular exponentiation, QFT, measurement
% - Continued fractions for period extraction
% - Extension to discrete logarithms and elliptic curves
% - Physics analogy: diffraction and interference (from notes)
% - Cite: Shor 1994/1997

\section{Grover's Algorithm and Amplitude Amplification}
\label{sec:grover}

% TODO: Adapt from notes.tex §3.1.2
% - Unstructured search problem
% - Grover iteration: oracle, diffusion
% - Optimality: BBBV lower bound
% - Impact on symmetric cryptography: AES-128 → AES-256
% - Amplitude amplification as generalization
% - Cite: Grover 1996, BBBV 1997

\section{Quantum Resource Estimates}
\label{sec:resource-estimates}

% TODO: Expand from notes.tex §3.1.1 (resource discussion)
% - Table: RSA-2048, RSA-4096, ECC-256 — logical qubits, T-gates, time
% - Physical qubit overhead with surface codes
% - Current hardware vs requirements (orders of magnitude gap)
% - Cite: Gidney-Eker\aa{} 2021, recent resource estimates

\section{What Quantum Computers Cannot Do}
\label{sec:quantum-limits}

% TODO: New section
% - NP-complete problems: no known exponential quantum speedup
% - Lattice problems: best quantum algorithms are polynomial speedup only
% - The polynomial/exponential divide as design criterion for PQC
% - Cite: Bennett et al. 1997, Regev 2004

%% Exercises
\begin{prob}
\label{prob:ch04-shor}
Walk through Shor's algorithm to factor $N = 15$ using the base $a = 7$. Show each step: modular exponentiation, QFT, measurement, continued fractions, and extraction of the factors.
\end{prob}

\begin{prob}
\label{prob:ch04-grover}
Grover's algorithm achieves a quadratic speedup for unstructured search. Explain why this means AES-128 offers only 64-bit security against a quantum adversary, while AES-256 retains 128-bit security. What is the implication for key size recommendations?
\end{prob}

\begin{prob}
\label{prob:ch04-resources}
Using the estimates from Gidney and Eker\aa{} (2021), calculate the number of physical qubits required to factor RSA-2048 assuming a surface code with physical error rate $10^{-3}$ and code distance $d = 27$. Compare this with the qubit counts of current quantum processors.
\end{prob}
